% from ega3-6-fr.tex
\oldpage[III]{154}
de $\mathfrak U^{(i)}$ sont affines (I, 5.5.6 et 1.3.11),
$\mathbf{Tor}_\bullet^S
(\mathfrak U^{(1)},\mathfrak U^{(2)};
 \mathscr P_\bullet^{(1)},\mathscr P_\bullet^{(2)})$ est un
bi-$\partial$-foncteur covariant en $\mathscr P_\bullet^{(1)}$,
$\mathscr P_\bullet^{(2)}$, à valeurs dans la catégorie des $A$-modules
($0_{\mathrm I}$, 11.7.3). En outre, on sait ($0_{\mathrm I}$, 11.7.2) que ce
bifoncteur est l'aboutissement commun de six foncteurs spectraux biréguliers,
que nous désignerons par la notation
\[
  {}^{(t)}\mathscr E^S
  (\mathfrak U^{(1)},\mathfrak U^{(2)};
   \mathscr P_\bullet^{(1)},\mathscr P_\bullet^{(2)})
\]
ou
\[
  {}^{(t)}\mathscr E
  (\mathfrak U^{(1)},\mathfrak U^{(2)};
   \mathscr P_\bullet^{(1)},\mathscr P_\bullet^{(2)}),
\]
où $t$ doit être remplacé par une des lettres $a$, $b$, $a'$, $b'$, $c$,
$d$, et dont les termes $E_2$ sont les suivants :
\[
\begin{aligned}
  {}^{(a)}\mathscr E_{pq}^2
  &=\bigoplus_{q_1+q_2=q}
    \mathbf{Tor}_p^A
    \bigl(\mathrm H_{q_1}(L_{\bullet\bullet}^{(1)}),
          \mathrm H_{q_2}(L_{\bullet\bullet}^{(2)})\bigr),\\
  {}^{(b)}\mathscr E_{pq}^2
  &=\mathrm H_p\bigl(
    \mathbf{Tor}_q^{A,\mathrm{II}}
    (L_{\bullet\bullet}^{(1)},L_{\bullet\bullet}^{(2)})\bigr),\\
  {}^{(a')}\mathscr E_{pq}^2
  &=\bigoplus_{q_1+q_2=q}
    \mathbf{Tor}_p^A
    \bigl(\mathrm H_{q_1}^{\mathrm I}(L_{\bullet\bullet}^{(1)}),
          \mathrm H_{q_2}^{\mathrm I}(L_{\bullet\bullet}^{(2)})\bigr),\\
  {}^{(b')}\mathscr E_{pq}^2
  &=\mathrm H_p\bigl(
    \mathbf{Tor}_q^A
    (L_{\bullet\bullet}^{(1)},L_{\bullet\bullet}^{(2)})\bigr),\\
  {}^{(c)}\mathscr E_{pq}^2
  &=\bigoplus_{q_1+q_2=q}
    \mathbf{Tor}_p^A
    \bigl(\mathrm H_{q_1}^{\mathrm{II}}(L_{\bullet\bullet}^{(1)}),
          \mathrm H_{q_2}^{\mathrm{II}}(L_{\bullet\bullet}^{(2)})\bigr),\\
  {}^{(d)}\mathscr E_{pq}^2
  &=\mathrm H_p\bigl(
    \mathbf{Tor}_q^{A,\mathrm I}
    (L_{\bullet\bullet}^{(1)},L_{\bullet\bullet}^{(2)})\bigr),
\end{aligned}
\]
où les notations sont conformes à celles de la théorie générale de
l'hyperhomologie. Nous allons dans ce qui suit expliciter davantage ces
termes initiaux.
\end{env}

\begin{env}[6.6.3]
\emph{Suites spectrales $(a)$ et $(a')$.} Nous avons vu en (6.6.1) que le
module d'homologie $\mathrm H_n(L_{\bullet\bullet}^{(i)})$ du bicomplexe
$L_{\bullet\bullet}^{(i)}$ était égal à
$\mathbf H^{-n}(X^{(i)},\mathscr P_\bullet^{(i)})$; donc
\[
  {}^{(a)}\mathscr E_{pq}^2
  =\bigoplus_{q_1+q_2=q}
   \mathbf{Tor}_p^A
   \bigl(\mathbf H^{-q_1}(X^{(1)},\mathscr P_\bullet^{(1)}),
         \mathbf H^{-q_2}(X^{(2)},\mathscr P_\bullet^{(2)})\bigr).
\]
Par définition, le complexe
$\mathrm H_n^{\mathrm I}(L_{\bullet\bullet}^{(i)})$ a pour terme de degré $k$
le module d'homologie
$\mathrm H_n(C^\bullet(\mathfrak U^{(i)},\mathscr P_k^{(i)}))$, c'est-à-dire,
par définition, le module de cohomologie
$\mathbf H^{-n}(\mathfrak U^{(i)},\mathscr P_k^{(i)})$; on sait (1.4.1) que
ce module est canoniquement isomorphe à
$\mathbf H^{-n}(X^{(i)},\mathscr P_k^{(i)})$; donc
\[
  {}^{(a')}\mathscr E_{pq}^2
  =\bigoplus_{q_1+q_2=q}
   \mathbf{Tor}_p^A
   \bigl(\mathbf H^{-q_1}(X^{(1)},\mathscr P_\bullet^{(1)}),
         \mathbf H^{-q_2}(X^{(2)},\mathscr P_\bullet^{(2)})\bigr).
\]
\end{env}

\begin{env}[6.6.4]
\emph{Suites spectrales $(b)$ et $(b')$.} Par définition,
$\mathbf{Tor}_q^{A,\mathrm{II}}
(L_{\bullet\bullet}^{(1)},L_{\bullet\bullet}^{(2)})$ est un bicomplexe dont
le terme de degré $(h,k)$ est le $A$-module
\[
  \mathbf{Tor}_q^A
  \bigl(C^{-h}(\mathfrak U^{(1)},\mathscr P_\bullet^{(1)}),
        C^{-k}(\mathfrak U^{(2)},\mathscr P_\bullet^{(2)})\bigr).
\]
Soit $\Phi^{(i)}$ l'ensemble d'indices de $\mathfrak U^{(i)}$; par
définition, le complexe de modules
$C^r(\mathfrak U^{(i)},\mathscr P_\bullet^{(i)})$ ($r\geqslant0$) est somme
directe des complexes
$\Gamma(U_\rho^{(i)},\mathscr P_\bullet^{(i)})$, où $U_\rho^{(i)}$ est
l'intersection des $U_\xi^{(i)}$ pour $\xi\in\rho$, et $\rho$ parcourt
$\mathfrak P(\Phi^{(i)})$; donc le $A$-module
\[
  \mathbf{Tor}_q^A
  \bigl(C^{-h}(\mathfrak U^{(1)},\mathscr P_\bullet^{(1)}),
        C^{-k}(\mathfrak U^{(2)},\mathscr P_\bullet^{(2)})\bigr)
\]
est somme directe des $A$-modules
\[
  \mathbf{Tor}_q^A
  \bigl(\Gamma(U_\sigma^{(1)},\mathscr P_\bullet^{(1)}),
        \Gamma(U_\tau^{(2)},\mathscr P_\bullet^{(2)})\bigr),
\]
où $\sigma$ (resp. $\tau$) parcourt les éléments de
$\mathfrak P(\Phi^{(1)})$ (resp. $\mathfrak P(\Phi^{(2)})$) tels que
$\operatorname{Card}(\sigma)=-(h+1)$ (resp.
$\operatorname{Card}(\tau)=-(k+1)$). Comme $X^{(1)}$ et $X^{(2)}$ sont des
schémas, les $U_\rho^{(i)}$ sont affines, donc on a (6.4.1.1)
\[
  \mathbf{Tor}_q^A
  \bigl(\Gamma(U_\sigma^{(1)},\mathscr P_\bullet^{(1)}),
        \Gamma(U_\tau^{(2)},\mathscr P_\bullet^{(2)})\bigr)
  =\Gamma\bigl(U_\sigma^{(1)}\times_S U_\tau^{(2)},
    \mathscr{T}\!or_q^S
    (\mathscr P_\bullet^{(1)},\mathscr P_\bullet^{(2)})\bigr).
\]
